% Generated draft: requires mathematical and source review.
\documentclass[11pt]{article}
\usepackage[margin=25mm]{geometry}
\usepackage{fontspec}
\setmainfont{Latin Modern Roman}
\usepackage{amsmath,amssymb,mathtools,unicode-math}
\setmathfont{Latin Modern Math}
\usepackage{enumitem,needspace,xurl,hyperref}
\setlength{\parindent}{0pt}
\setlength{\parskip}{5pt}
\setlength{\emergencystretch}{3em}
\newcommand{\N}{\mathbb{N}}
\newcommand{\Z}{\mathbb{Z}}
\newcommand{\Q}{\mathbb{Q}}
\newcommand{\R}{\mathbb{R}}
\newcommand{\C}{\mathbb{C}}
\newcommand{\F}{\mathbb{F}}
\newcommand{\A}{\mathbb{A}}
\newcommand{\PP}{\mathbb{P}}
\newcommand{\E}{\mathbb{E}}
\newcommand{\eps}{\varepsilon}
\newcommand{\dd}{\,\mathrm{d}}
\begin{document}
\section*{UnsolvedMath ID: 1101310}
\textbf{Source problem code:} catalogue entry
\subsection{Definitions and mathematical statement}
Let $X$ be a finite set of generators and write $X^{\pm1}=X\cup X^{-1}$. Let $W=w_1\cdots w_m$ be a cyclically reduced word: there is no adjacent inverse pair, including the pair of its final and initial letters when $m>0$. Set
\[
 G=\langle X\mid W=1\rangle,
 \qquad q_i=w_1\cdots w_i\quad(0\leq i\leq m),
\]
where $q_0$ is the empty word. Denote evaluation in $G$ by $[\ \cdot\ ]_G$. The prefix submonoid is
\[
 P_W=\left\{[q_{i_1}\cdots q_{i_k}]_G:
 k\geq0,\quad 0\leq i_j\leq m\right\}.
\]
Thus only multiplication and the identity are required; inverses of prefixes are not automatically included. The empty and full prefixes represent the identity and do not change the generated monoid.

The question is whether membership in $P_W$ is decidable for every such presentation: given a word $u\in(X^{\pm1})^*$, is there an algorithm that halts and decides whether $[u]_G\in P_W$? In the per-presentation reading this means
\[
 \forall (X,W)\ \exists A_{X,W}\ \forall u,
 \quad A_{X,W}(u)\text{ halts and correctly decides }[u]_G\in P_W.
\]
A uniform version instead requires one algorithm accepting $(X,W,u)$ as input and valid for all cyclically reduced $W$. The source question does not expressly distinguish these quantifier orders, while the supplied assessment discusses absence of a uniform algorithm. The finite generating alphabet is the standard effective presentation setting; the isolated source does not explicitly state its finiteness.
\subsection{Short English statement}
For a one-relator group, can one decide whether an element is a product of prefixes of its cyclically reduced defining relator? Cyclic reduction is an essential hypothesis.

The supplied assessment labels this case open. It reports negative answers for merely reduced relators, citing Gray (2020) and Foniqi--Gray--Nyberg-Brodda (2025), the latter described there as using a quasi-positive freely reduced relator $uv^{-1}$. It reports substantial decidable cyclically reduced families, citing Margolis--Meakin--Šuniḱ (2005), Juhász (2014), and Dolinka--Gray (2021), but no known cyclically reduced counterexample or uniform algorithm. It also reports an equivalence, via Ivanov--Margolis--Meakin (2001), with the word problem for cyclically reduced one-relator inverse monoids, and says a positive answer would imply decidability for all one-relation monoids. The exact algorithmic quantifiers and hypotheses of these reductions are not supplied; the reports are not new results or independently verified current-status claims.
\subsection{Sources}
\begin{itemize}
\item UnsolvedMath permanent ID 1101310, source code catalogue entry, ``Questions in Geometric Group Theory --- Q 13.10.'' Registry curation and metadata: Ulam AI and UnsolvedMath Contributors, CC BY 4.0; updated 2026-08-06.
\item Bestvina, \emph{Questions in Geometric Group Theory} (source list dated 2004), Question 13.10, PDF page 25: \url{https://www.math.utah.edu/~bestvina/eprints/questions-updated.pdf}. Supplied manual transcription accessed 2026-07-29. Original metadata records partial progress or special cases, with current-status release review and rights review still pending.
\item Gray (2020); Foniqi--Gray--Nyberg-Brodda (2025); Margolis--Meakin--Šuniḱ (2005); Juhász (2014); Dolinka--Gray (2021); Ivanov--Margolis--Meakin (2001), as named in the supplied OPEN-TRIAGE assessment. No titles, venues, theorem numbers, or locators are supplied for these references; none are inferred. The spelling Šuniḱ is retained from the record.
\end{itemize}
\end{document}
